\chapter{Joint Continuous RVs, Marginal Densities \& Conditional Densities}

Extending continuous distributions to multivariate spaces requires calculus over regions in $\mathbb{R}^2$. In machine learning, continuous multivariate densities govern everything from multi-sensor feature spaces to bivariate Gaussian classification. In this chapter, we master \textbf{Joint PDFs $f_{X,Y}(x, y)$}, \index{Marginal Densities}\textbf{Marginal Densities}, the geometric requirements for \index{Continuous Independence}\textbf{Continuous Independence}, \textbf{Conditional Densities $f_{X|Y}(x \mid y)$}, and \index{Conditional Expectation}\textbf{Conditional Expectation} $\E[Y \mid X]$ as an optimal predictor.

\vspace{0.5cm}
\begin{tcolorbox}[colback=primarylight, colframe=primarydark, title={\textbf{What You'll Learn in this Chapter}}]
\begin{itemize}
    \item Double integration and normalization axioms of bivariate continuous density functions.
    \item Integrating out nuisance variables over complex and variable-dependent boundaries.
    \item The Support Rectangle Theorem: why non-rectangular supports strictly forbid independence.
    \item Slicing joint surfaces into normalized conditional density curves $f_{X|Y}(x \mid y)$.
    \item The Law of Iterated Expectations (Adam's Law): $\E[\E[Y \mid X]] = \E[Y]$.
\end{itemize}
\end{tcolorbox}

\section{Joint Probability Density Function}

\begin{definitionbox}{Joint Continuous PDF $f_{X,Y}(x, y)$}
Two continuous random variables $X$ and $Y$ possess a \index{Joint PDF}\textbf{Joint PDF} $f_{X,Y}(x, y) \ge 0$ such that the probability of the pair $(X, Y)$ landing within any two-dimensional Borel region $\mathcal{A} \subseteq \mathbb{R}^2$ is the volume under the density surface:
$$\P((X, Y) \in \mathcal{A}) = \iint_{\mathcal{A}} f_{X,Y}(x, y) \, dx \, dy$$
\textbf{Total Volume Normalization Axiom:}
$$\int_{-\infty}^{\infty} \int_{-\infty}^{\infty} f_{X,Y}(x, y) \, dx \, dy = 1$$
\end{definitionbox}

\begin{videocard}{IITM BS Lecture: Joint Continuous Random Variables}{https://www.youtube.com/watch?v=CIyqDJG-U9Q}
Official lecture introducing 2D density surfaces, double integrals over regions, and normalization.
\end{videocard}

\section{Marginal Densities and Independence}

\begin{intuitionbox}{Marginalization as a Geometric Silhouette}
Visualize a 2D continuous joint PDF as a mountain range sitting on the $(x, y)$ plane. The Marginal PDF $f_X(x)$ represents the shadow or silhouette cast on the $x$-axis if you project light parallel to the $y$-axis. By integrating out $y$, you collapse all probability mass along the cross-sectional slice into a single coordinate point on the margin!
\end{intuitionbox}

\begin{lstlisting}[language=Python, caption=Data Science: Numerical Double Integration in SciPy]
import numpy as np
from scipy import integrate

# Joint PDF: f(x, y) = 2 over triangular domain 0 <= y <= x <= 1
def joint_pdf(y, x):
    return 2.0 if (0 <= y <= x <= 1) else 0.0

# Marginal f_X(x) for fixed x = 0.6: integrate y from 0 to x
x_val = 0.6
marginal_x, _ = integrate.quad(lambda y: joint_pdf(y, x_val), 0, x_val)
print(f"Marginal f_X(0.6) = {marginal_x:.2f}") # Expected: 2 * 0.6 = 1.2
\end{lstlisting}

\begin{formulabox}{Marginal Density Formulas}
To obtain the univariate density of $X$ (or $Y$), integrate over the entire range of the other variable:
\begin{itemize}
    \item \textbf{Marginal PDF of $X$:} $f_X(x) = \int_{-\infty}^{\infty} f_{X,Y}(x, y) \, dy$
    \item \textbf{Marginal PDF of $Y$:} $f_Y(y) = \int_{-\infty}^{\infty} f_{X,Y}(x, y) \, dx$
\end{itemize}
\end{formulabox}

\begin{definitionbox}{Continuous Independence Criterion}
Continuous random variables $X$ and $Y$ are \index{Independent}\textbf{Independent} ($X \perp\!\!\!\perp Y$) if and only if for all $(x, y) \in \mathbb{R}^2$:
$$f_{X,Y}(x, y) = f_X(x) \cdot f_Y(y)$$
\textbf{Necessary Geometric Prerequisite:} The joint support $\mathcal{S}_{X,Y} = \{(x, y) \colon f_{X,Y}(x, y) > 0\}$ must be a \textbf{product rectangle} $\mathcal{S}_X \times \mathcal{S}_Y$ (which may be infinite or semi-infinite). If the boundary of the support involves both variables (e.g., $y \le x$ or $x^2 + y^2 \le 1$), the variables are \textbf{never independent}!
\end{definitionbox}

\begin{videocard}{IITM BS Lecture: Marginal Densities of Continuous Random Variables}{https://www.youtube.com/watch?v=TirRvH4PiiA}
Official lecture deriving marginal PDF slices via partial integration.
\end{videocard}

\begin{videocard}{IITM BS Lecture: Independence of Multiple Continuous Random Variables}{https://www.youtube.com/watch?v=cms_2NGMiHY}
Official lecture testing product factorization for joint continuous density functions.
\end{videocard}

\section{Conditional Densities and Conditional Expectation}

\begin{definitionbox}{Conditional PDF $f_{Y|X}(y \mid x)$}
For any fixed $x$ where $f_X(x) > 0$, the \index{Conditional Density}\textbf{Conditional PDF} of $Y$ given $X = x$ is:
$$f_{Y|X}(y \mid x) = \frac{f_{X,Y}(x, y)}{f_X(x)}$$
Integrating over all $y$ yields $\int_{-\infty}^{\infty} f_{Y|X}(y \mid x) \, dy = \frac{\int f_{X,Y}(x,y)dy}{f_X(x)} = \frac{f_X(x)}{f_X(x)} = 1$.
\end{definitionbox}

\begin{theorembox}{The Law of Iterated Expectations (Adam's Law)}
\index{Law of Iterated Expectations}\index{Adam's Law}
The conditional expectation $\E[Y \mid X]$ is a function of the random variable $X$, and is therefore itself a \textbf{random variable}. Its expected value over the distribution of $X$ recovers the unconditional mean of $Y$:
$$\E\big[ \E[Y \mid X] \big] = \E[Y]$$
\textbf{Proof:}
$$\E\big[ \E[Y \mid X] \big] = \int_{-\infty}^{\infty} \left( \int_{-\infty}^{\infty} y f_{Y|X}(y \mid x) \, dy \right) f_X(x) \, dx = \int_{-\infty}^{\infty} y \left( \int_{-\infty}^{\infty} f_{X,Y}(x, y) \, dx \right) dy = \int_{-\infty}^{\infty} y f_Y(y) \, dy = \E[Y]$$
\end{theorembox}

\begin{videocard}{IITM BS Lecture: Conditional Density of Continuous Random Variables}{https://www.youtube.com/watch?v=lTCc9oYRwr8}
Official lecture constructing conditional density curves and computing conditional expectations.
\end{videocard}

\section{Exam Gotchas \& Cognitive Traps}

\begin{trapbox}{The Triangular Limit Integration Trap}
\textbf{The Trap:} When finding the marginal density $f_X(x)$ over a dependent region such as $0 \le y \le x \le 1$, integrating with static constant limits: $f_X(x) = \int_0^1 f(x, y) \, dy$.

\textbf{The Reality:} You \textbf{must trace the active boundary for the specific variable being integrated}:
\begin{itemize}
    \item To find $f_X(x)$, fix $x \in [0, 1]$. In the inequality $0 \le y \le x \le 1$, $y$ is constrained between $0$ and $x$.
    $$f_X(x) = \int_0^x f_{X,Y}(x, y) \, dy$$
    \item To find $f_Y(y)$, fix $y \in [0, 1]$. In the inequality $0 \le y \le x \le 1$, $x$ is constrained between $y$ and $1$.
    $$f_Y(y) = \int_y^1 f_{X,Y}(x, y) \, dx$$
\end{itemize}
Writing constant limits on coupled boundaries is the single most common deduction in calculus-based probability exams!
\end{trapbox}

\begin{trapbox}{The Product Factorization Illusion on Non-Rectangles}
\textbf{The Trap:} Looking at $f_{X,Y}(x, y) = 8xy$ on $0 \le x \le y \le 1$ and concluding that $X \perp\!\!\!\perp Y$ because $8xy$ factorizes algebraically as $(8x) \cdot (y)$.

\textbf{The Reality:} Independence requires that $f_{X,Y}(x, y) = f_X(x) f_Y(y)$ \textbf{everywhere on $\mathbb{R}^2$}, including the indicator function of the support:
$$f_{X,Y}(x, y) = 8xy \cdot \mathbb{I}(0 \le x \le y \le 1)$$
The indicator $\mathbb{I}(0 \le x \le y \le 1)$ cannot be factorized into $\mathbb{I}_X(x) \cdot \mathbb{I}_Y(y)$! The point $(x=0.8, y=0.2)$ has $f_{X,Y}(0.8, 0.2) = 0$, but both marginals are positive: $f_X(0.8) > 0$ and $f_Y(0.2) > 0$. Therefore, $X$ and $Y$ are \textbf{dependent}.
\end{trapbox}

\section{Comparative Synthesis Matrix}

\begin{table}[htbp]
\centering
\small
\begin{tabular}{p{3.2cm}p{3.8cm}p{3.8cm}p{3.5cm}}
\toprule
\textbf{Support Geometry} & \textbf{Rectangular Domain} & \textbf{Triangular Domain} & \textbf{Unit Disk Domain} \\
\midrule
\textbf{Domain Definition} & $[a, b] \times [c, d]$ (Independent limits) & $0 \le y \le x \le 1$ (Coupled limits) & $x^2 + y^2 \le r^2$ (Quadratic coupling) \\
\addlinespace
\textbf{Marginal $f_X(x)$ Limits} & $\int_c^d f(x, y) \, dy$ (Constant) & $\int_0^x f(x, y) \, dy$ (Function of $x$!) & $\int_{-\sqrt{r^2-x^2}}^{\sqrt{r^2-x^2}} f(x, y) \, dy$ \\
\addlinespace
\textbf{Can RVs be Independent?} & \textbf{Yes}, if algebraic formula factorizes & \textbf{Never} (Forbidden by boundary coupling) & \textbf{Never} (Forbidden by circular boundary) \\
\addlinespace
\textbf{Conditional $f_{Y|X}(y|x)$ Support} & Constant interval $[c, d]$ & Variable interval $[0, x]$ & Variable interval $[-\sqrt{r^2-x^2}, \sqrt{r^2-x^2}]$ \\
\addlinespace
\textbf{Data Science Analog} & Uncorrelated orthogonal bounding boxes & Hierarchical pipeline latencies (Step 2 requires Step 1) & 2D spatial / GPS error dispersion \\
\bottomrule
\end{tabular}
\caption{Classification of Bivariate Continuous Integration Regimes by Support Geometry.}
\end{table}

\newpage
\section{Practice Exercises}
\textit{Detailed step-by-step solutions are available in the Appendix.}

\subsection*{Part A: Analytical \& Computational Problems}

\begin{enumerate}
\item \textbf{Joint Density Normalization Constant:} Let $X$ and $Y$ have joint PDF:
    $$f_{X,Y}(x, y) = \begin{cases} c x y & \text{for } 0 \le x \le 1, \, 0 \le y \le 2 \\ 0 & \text{otherwise} \end{cases}$$
    Find constant $c$ and compute $\P(X \le 0.5, Y \le 1)$.

\item \textbf{Marginal Densities Computation:} For the joint PDF $f_{X,Y}(x, y) = c xy$ on $0 \le x \le 1, 0 \le y \le 2$ from Exercise 1:
    \begin{enumerate}
        \item Compute marginal density $f_X(x)$.
        \item Compute marginal density $f_Y(y)$.
        \item Are $X$ and $Y$ independent?
    \end{enumerate}

\item \textbf{Triangular Region Joint PDF:} Let $X$ and $Y$ be uniformly distributed over the triangle $0 \le y \le x \le 1$ with joint PDF $f_{X,Y}(x, y) = 2$ for $0 \le y \le x \le 1$.
    Find marginal density $f_X(x)$ and $f_Y(y)$.

\item \textbf{Conditional Density Calculation:} For the triangular joint PDF in Exercise 3:
    Find conditional density $f_{Y|X}(y \mid x)$ and calculate conditional expectation $\E[Y \mid X = x]$.

\item \textbf{Data Science Application (Bivariate Normal Feature Regression):} Standardized features $(X, Y)$ follow a Bivariate Normal distribution with $\mu_X = 0, \mu_Y = 0, \sigma_X = 1, \sigma_Y = 1$ and correlation $\rho = 0.8$. Given the formula $\E[Y \mid X = x] = \mu_Y + \rho \frac{\sigma_Y}{\sigma_X} (x - \mu_X)$, predict the expected value of $Y$ when feature $X = 2.5$.
\end{enumerate}

\subsection*{Part B: Conceptual Diagnostic Traps (Exam Gotchas)}

\begin{enumerate}[resume]
\item \textbf{Diagnostic Trap (Support Boundary Independence Criterion):}
Let $X$ and $Y$ have joint density $f_{X,Y}(x, y) = 8xy$ defined over the triangular support $\mathcal{T} = \{(x, y) \colon 0 \le x \le y \le 1\}$. A student claims $X$ and $Y$ are independent because $8xy = (8x) \cdot (y)$. Is the student correct?
\begin{enumerate}
    \item Yes, algebraic factorization of the density expression guarantees statistical independence.
    \item Yes, because the marginal densities $f_X(x)$ and $f_Y(y)$ multiply to $8xy$.
    \item No, because the joint support is non-rectangular; the range of $X$ depends on $Y$, preventing independence.
    \item No, because independence can only occur when the constant coefficient equals 1.
\end{enumerate}

\item \textbf{Diagnostic Trap (The Nature of Conditional Expectation):}
Which of the following statements regarding the conditional expectation $\E[Y \mid X]$ is strictly correct?
\begin{enumerate}
    \item $\E[Y \mid X]$ is a constant fixed number equal to $\E[Y]$.
    \item $\E[Y \mid X]$ is a random variable that is a function of $X$, and its expected value is $\E[Y]$.
    \item $\E[Y \mid X] = \frac{\E[XY]}{\E[X]}$.
    \item $\E[Y \mid X]$ is only defined if $X$ and $Y$ are jointly Gaussian.
\end{enumerate}
\end{enumerate}
